% !TEX root = ../ApproxDA-TransUNet.tex
\section{Introduction}

Transformers have advanced medical image segmentation by modeling long-range dependencies that convolutional networks capture only indirectly. Dual attention~\cite{fu2019dual}---a Position Attention Module (PAM) for spatial correlations and a Channel Attention Module (CAM) for channel correlations---is one such mechanism, and DA-TransUNet~\cite{sun2024datransunet} showed its effectiveness in a hybrid CNN--Transformer architecture for multi-organ and polyp segmentation. Full PAM costs $\mathcal{O}(N^2C)$ for $N$ positions and $C$ channels, and windowed attention~\cite{liu2021swin}, low-rank projection~\cite{wang2020linformer}, and grouped channel interaction~\cite{rahman2024emcad} are standard ways to reduce this cost. These approximations are usually adopted as accuracy-neutral engineering choices, which raises a question that is rarely examined across tasks with different spatial characteristics:

\begin{center}
\emph{Does attention approximation change segmentation accuracy,\\
and does the effect depend on the task?}
\end{center}

To study this question we design \textbf{ApproxDA-TransUNet}\footnote{Source code will be made available upon acceptance.}, which keeps the architecture of DA-TransUNet and replaces only its attention modules with approximations controlled by three independent parameters. With the backbone, training protocol, and evaluation fixed, we vary one parameter at a time on three datasets with different anatomical and visual characteristics, and compare against DA-TransUNet re-trained under the same protocol.

The main contributions of this paper are:
\begin{itemize}
\item \textbf{A controlled dual-attention approximation framework.} Window size $M$, rank $r$, and group number $G$ control the spatial scope, the affinity resolution, and the channel-interaction scope of dual attention. In its exact limit the block reduces to the DA-TransUNet block, so each approximation can be ablated in isolation.
\item \textbf{Evidence of dataset-dependent window sensitivity.} Synapse multi-organ CT is substantially more sensitive to $M$ than Kvasir-SEG and ISIC~2018, and with an appropriately selected window ApproxDA-TransUNet improves Dice over the matched DA-TransUNet on all three datasets.
\end{itemize}

The remainder of this paper is organized as follows: Section~\ref{sec:related_work} reviews related work; Section~\ref{sec:method} presents ApproxDA-TransUNet; Section~\ref{sec:experiments} describes the protocol and main results; Section~\ref{sec:understanding} reports the controlled ablations; Section~\ref{sec:conclusion} concludes.
